\begin{annexitem}{annex:p12.har.prcompiler-typing}{thm:p12-har-prcompiler-typing}{Typing the complete finite prescription compiler}
\begin{proof}
Induct on the finite ranked prescription, with its arity as part of the
induction assertion. The type \(T_0=\mathsf{nat}\) is admitted; if
\(T_j:\mathsf{Set}\), the admitted Set-to-Set product gives
\(\mathsf{nat}\to T_j:\mathsf{Set}\). Thus every \(T_k\) is
an admitted data type. In the context
\(x_1:\mathsf{nat},\ldots,x_k:\mathsf{nat}\), O:nat gives the
zero template after k abstractions. For k=0 it gives O itself, not a
function on a fictitious empty-tuple value. Sx:nat gives the unary
successor template. The side condition \(1\le i\le k\) makes
\(x_i:\mathsf{nat}\) a declared variable, so projection has type
\(T_k\); there is deliberately no zeroary projection case.

For composition, the induction hypotheses give
\(\mathcal C f:T_m\) and \(\mathcal C g_i:T_k\). Weakening
these closed templates into the x context and k successive applications
give \((\mathcal C g_i)\vec x:\mathsf{nat}\). Apply
\(\mathcal C f\) to these m arguments, one at a time, to obtain nat.
Then k abstractions give \(T_k\). At m=0 this is the judgment
\(\mathcal C f:T_0=\mathsf{nat}\) with no applications. The
unused x declarations are simply weakened declarations.

For recursion, the induction hypotheses give \(\mathcal C g:T_k\)
and \(\mathcal C h:T_{k+2}\). In the x context the seed
\((\mathcal C g)\vec x\) has type nat. In its extension by
\(n:\mathsf{nat},r:\mathsf{nat}\), exactly k+2 applications give
\((\mathcal C h)\vec x\ n\ r:\mathsf{nat}\). Abstract over r
then n to get the step type nat-to-nat-to-nat. The fixed-target nat
recursor therefore returns nat-to-nat; k abstractions give
\(T_{k+1}\). At k=0 the same judgments use g:nat and
h:nat-to-nat-to-nat, with no x declarations. None of these steps reads
an existential proof to choose the seed or the next value.

The construction traverses a finite prescription once to emit each
template and its subtemplates, so it is an actual metalevel algorithm.
No normalization or classifier of arbitrary source terms is called.
\end{proof}
\end{annexitem}

\begin{annexitem}{annex:p12.har.prcompiler-predicate}{thm:p12-har-prcompiler-predicate}{The explicit arbitrary-prescription characteristic program and computation}
\begin{proof}
Typing the compiled \(e\) in the \(k\)-variable nat context gives a nat
result; \(\mathsf{one}:\mathsf{nat}\) then makes its equality constructor a form value.
Exactly \(k\) abstractions give \(p_e\), and the predicate-instance theorem
gives the typed nat-valued \(K_e\) and its range. Expanding the
\(\mathcal P\) interpretation of this form is exactly the internal
equation \((\mathcal C e)\vec x=\mathsf{one}\). Thus the two
internal iff assertions follow from the positive and negative maps
already constructed, without assuming the numerical predicate's
conclusion. They apply to arbitrary typed open variables.

For the additional numerical computation statement, fix a numerical
tuple and put \(v:=\langle e\rangle(\vec a)\). The actual
representation reduction supplies
\((\mathcal C e)\overline{\vec a}\longrightarrow^*\overline v\).
Unfold \(p_e\), \(K_e\) and the atomic comp clause. Congruence permits that
finite reduction in the first \(\mathsf{b\_eq}\) argument, leaving
\(\mathsf{code}(\mathsf{b\_eq}\ \overline v\ (S O))\).
If \(v=0\), the zero/successor equality clause gives false, then code
gives \(O\). If \(v=1\), successor/successor followed by zero/zero gives
true, then code gives one. If \(v\) is at least two, write its numeral
as \(S(S\overline w)\); the successor/successor clause leaves
b\_eq\((S\overline w)O\), whose successor/zero clause gives
false, then code gives O. Ordinary natural-number case analysis
has exhausted every possible v, so exactly v=1 yields one and
every other v yields O. The k=0 instance has the same calculation
with no input applications. These are prescribed finite contractions,
not a normalization theorem for the entire host calculus.

A predicate given by a finite zero/one characteristic prescription
has truth definition precisely value one, so the preceding construction
represents that original predicate. No new assumption of a desired
adequacy theorem was introduced. The output program's range was
derived from its boolean data selector, independently of any claimed
range for the input prescription. Thus the same construction is
well defined for every finite prescription, while an arbitrary host
predicate is not declared primitive recursive without such a prescription.
\end{proof}
\end{annexitem}

\begin{annexitem}{annex:p12.har.prcompiler-representation}{thm:p12-har-prcompiler-representation}{The complete universal numerical representation derivation}
\begin{proof}
First ordinary induction on n specifies the numerical value in the
recursion clause: at zero use the already specified value of g; at
n+1 use the already specified value at n as h's last argument.
Induction on the finite prescription therefore defines a total numerical
function at its exact arity. This does not use the emitted term to
define its purported meaning.

Now induct on the prescription, universally quantifying the numerical
argument tuple. The zero template beta-contracts to O, which is the
numeral for zero, including the no-abstraction k=0 case. The successor
template on \(\overline a\) beta-contracts to
\(S\overline a=\overline{a+1}\). The projection template contracts
to its ith input numeral. For composition, the k outer beta steps
expose the application of the compiled f to the m compiled children.
By the child induction hypotheses, reduce the ith child in its argument
position to \(\overline{\langle g_i\rangle(\vec a)}\), for every
i from 1 through m. Congruence of reduction licenses these finite
argument-position reductions. The hypothesis for f then reduces the
whole application to the numeral for the independently defined
composition value. At m=0 no child reduction is performed and the
zeroary f hypothesis is applied directly.

In the recursion case fix \(\vec a\), and make a second induction
on the last numerical input n. Put
\(B:=(\mathcal C g)\overline{\vec a}\) and
\(H:=\lambda j:\mathsf{nat}.\lambda r:\mathsf{nat}.
(\mathcal C h)\overline{\vec a}\ j\ r\).
At zero the recursor contracts to B, and the g hypothesis gives the
required numeral. At n+1 the successor computation gives
\[
 \mathsf{nat\_rec}\ B\ H\ (S\overline n)
 \longrightarrow H\ \overline n\
                      (\mathsf{nat\_rec}\ B\ H\ \overline n).
\]
After the two beta steps, reduce the displayed recursive argument by
the inner induction hypothesis to
\(\overline{\langle\mathsf{Rec}_k(g,h)\rangle(\vec a,n)}\).
The hypothesis for h, at the k+2 now numerical arguments, gives the
numeral for precisely the successor clause of the independent meaning.
No simultaneous normalization of arbitrary open inputs was assumed:
these are explicit finite reduction sequences for every numerical tuple.

A valuation assigning ordinary natural numbers to the k parameters
supplies exactly such a tuple. Its numeral substitution is the displayed
calculation, so the result holds for every standard valuation. This
valuation statement is not an assertion of a chosen standard-model
interpretation for every uninterpreted carrier, nor of contraction of
an opaque variable to a numeral.
\end{proof}
\end{annexitem}

\begin{annexitem}{annex:p12.har.prcompiler-boolean}{thm:p12-har-prcompiler-boolean}{All boolean cases, types and the nonidentical code bridge}
\begin{proof}
The generated selector at bool takes a bool discriminator and two bool
branches; all four connective terms therefore have the displayed unary
or binary bool types. The selector at nat gives code:bool-to-nat.
Case analysis into the proposition code Prop gives istrue:bool-to-Prop;
this is elimination of a data constructor, not of a Prop proof into Set.
For the comparisons, \(q_0,e_0:\mathsf{nat}\to\mathsf{bool}\).
The inner recursion in d has target bool, false seed, and step
\(\lambda j:\mathsf{nat}.\lambda r:\mathsf{bool}.q j\) of
type nat-to-bool-to-bool. Thus d has type
\(\mathsf{nat}\to(\mathsf{nat}\to\mathsf{bool})\to
(\mathsf{nat}\to\mathsf{bool})\). The two outer recursions
have fixed target nat-to-bool and give the required binary comparison
types. The three derived comparisons apply these terms to nat arguments.
Expanding their recursors gives exactly the four constructor equations
stated in the body, with no assumed comparison primitive.

Here are full propositional constructions for the boolean connective
cases. Put T:=TT and F:=FF. The equivalences
\(T\wedge T\leftrightarrow T\), \(T\vee A\leftrightarrow T\)
and \(A\vee T\leftrightarrow T\) project or return trivial in the
forward direction and introduce the corresponding conjunction or
disjunction in the reverse direction. For
\(F\wedge A\leftrightarrow F\), project the F component and use
FF elimination for the reverse direction; the symmetric conjunction
is identical with the other projection. For
\(F\vee A\leftrightarrow A\), eliminate the disjunction, with
the F branch eliminated and the A branch returned; reverse by right
injection. The symmetric disjunction uses left injection. These maps
also settle F-with-F. Finally
\(\neg T\leftrightarrow F\) applies the alleged negation to
trivial, with FF elimination in reverse; \(\neg F\leftrightarrow T\)
returns trivial forward and the identity F-to-F map reverse.

For implication the four pairs (a,b) reduce the left side to,
respectively for (true,true), (true,false), (false,true), (false,false),
T,F,T,T. In the first case \(T\to T\leftrightarrow T\): forward
return trivial, reverse return the constant trivial function. In the
second \(T\to F\leftrightarrow F\): forward apply to trivial,
reverse eliminate F to construct the function. In the last two,
\(F\to T\leftrightarrow T\) and \(F\to F\leftrightarrow T\):
forward return trivial; reverse use FF elimination or the identity
F-to-F function. Combining these maps with the literal selector
computations proves implication. The preceding conjunction,
disjunction and negation maps prove the other connective equivalences
in all their literal cases. Nested generated bool induction supplies
these results for arbitrary a and b, without excluded middle.

For the code bridge, bool induction has only two branches. At true,
code reduces to one and istrue to T. Return the right range injection;
T-to-(one=one) returns equality reflexivity, and the reverse returns
trivial. Both \(\neg T\) and \(\mathsf{one}=O\) imply F:
the former by application to trivial, the latter by nat constructor
disjointness from the historical equality lemma. Conversely F
elimination gives either. Thus their equivalence is constructed. At
false, code reduces to O and istrue to F. Return the left range
injection; F and O=one are equivalent by FF elimination and nat
constructor disjointness. Both \(\neg F\) and O=O are inhabited:
use the identity F-to-F map and reflexivity to give both directions.
These branches prove all three assertions for arbitrary b. In
particular bool values were not mistaken for nat values.
\end{proof}
\end{annexitem}

\begin{annexitem}{annex:p12.har.prcompiler-atoms}{thm:p12-har-prcompiler-atoms}{The complete five-atom comparison proofs}
\begin{proof}
For b\_le, induct on n with motive
\(\Pi m:\mathsf{nat}.
\mathsf{istrue}(\mathsf{b\_le}\ n\ m)\leftrightarrow n\le m\).
At n=O the program returns true for every m. The historical order
lemma gives \(O\le m\); T-to-order returns that proof and
order-to-T returns trivial. At n=S a, perform nat case analysis on m.
At m=O the program returns false. No positive natural is at most O,
by the already proved historical order lemma; the order-to-F map is
that refutation and the F-to-order map is FF elimination. At m=S b
the program contracts to b\_le a b. The induction hypothesis gives
its istrue equivalence to \(a\le b\). Historical successor
monotonicity and cancellation give the two maps between
\(a\le b\) and \(S a\le S b\). Composing gives both maps
required in the successor/successor branch. These three branches
are the complete generalized induction.

For b\_eq use the same motive with equality. At n=O split m:
zero/zero returns true and both directions use equality reflexivity
or trivial; zero/successor returns false and the equality-to-F map
is constructor disjointness. At n=S a split m: successor/zero is
again the disjointness case, and successor/successor contracts to
b\_eq a b. The induction hypothesis supplies equality of a and b
in one direction, which successor congruence turns into equality of
their successors; in the reverse direction successor injectivity
supplies a=b, then the reverse induction map proves istrue. F
elimination supplies the other map in each disjointness case.
All equality manipulations have Prop targets, as established in
\cref{thm:p12-har-prime-equality}.

Finally b\_lt n m is literally b\_le(S n)m, and the historical
definition of n<m is \(S n\le m\). Apply the weak-order result at
these arguments. The b\_ge and b\_gt definitions reverse the
arguments, giving \(m\le n\) and m<n, respectively. Thus all five
source atomic interpretations, not only equality, have both maps.
\end{proof}
\end{annexitem}

\begin{annexitem}{annex:p12.har.prcompiler-bounds}{thm:p12-har-prcompiler-bounds}{The complete exclusive-bound induction and both directions}
\begin{proof}
The two seeds have type C and each step has type nat-to-C-to-C:
for f:nat-to-bool, f m and r f are bool, and their conjunction or
disjunction is bool. The recursors therefore have type nat-to-C,
giving the announced quantifier program types and literal zero and
successor computations.

We first establish the exact endpoint decomposition from the historical
order library, rather than assuming an enumeration theorem:
\[
 x<S m\ \leftrightarrow\ (x<m)\vee(x=m).
\]
Indeed \(x<S m\) means \(S x\le S m\); cancellation gives
\(x\le m\). Trichotomy for x,m gives x<m, x=m or m<x.
The last is impossible by x\(\le\)m and the proved strict-order
incompatibility, so the first two injections remain. Conversely, x<m
implies \(x\le m\), then \(S x\le S m\); x=m gives that same
inequality by equality transport from reflexivity. Thus both maps are
constructive. The order library also gives no x<O and \(m<S m\).

Induct on b, universally quantifying the data family f. For all at
zero, istrue of the seed is T. Its forward map sends any x and
x<O proof to FF elimination into istrue(f x). Its reverse map
returns trivial. For ex at zero, istrue of the seed is F.
The forward map eliminates F to the existential. The reverse map
eliminates a purported witness x and its x<O proof to F. Witness
elimination is into a proposition only.

For all at S m, constructor computation and the boolean conjunction
lemma identify istrue with
\(\mathsf{istrue}(f m)\wedge
 \mathsf{istrue}(\mathsf{b\_all}\ m\ f)\).
In the forward direction take this pair and apply the induction
hypothesis to its second component to get the old universal prefix.
For arbitrary x<S m use the endpoint decomposition. In the x<m
case apply that prefix; in the x=m case transport the first component
along the equality. This constructs the new universal family.
In the reverse direction, apply the proposed universal family at m
and \(m<S m\) for the endpoint value. Restrict the family to x<m
by the forward inclusion x<m implies x<S m; the reverse induction
map gives the old b\_all istrue proof. Pair the two results and apply
the reverse boolean-conjunction map. These are actual lambda, pair,
case and transport constructions for both implications.

For ex at S m, the analogous boolean disjunction identifies istrue
with \(\mathsf{istrue}(f m)\vee
\mathsf{istrue}(\mathsf{b\_ex}\ m\ f)\).
In its left branch produce existential witness m with \(m<S m\)
and the supplied value proof. In its right branch apply the forward
induction map, eliminate its existential only into the new existential
proposition, and use x<m implies x<S m on the same witness.
In reverse, eliminate a supplied witness x with x<S m and istrue(f x)
into the disjunction goal. Endpoint decomposition gives x<m or x=m.
The first builds a witness for the old existential, applies the
reverse induction map, and returns the right disjunction injection.
The second transports the value proof to f m and returns the left
injection. Then apply the reverse boolean-disjunction map. No Set
datum in either quantifier program depends on a Prop witness or on DN.
\end{proof}
\end{annexitem}

\begin{annexitem}{annex:p12.har.prcompiler-form-typing}{thm:p12-har-prcompiler-form-typing}{Typing both internal interpreters and all higher-order branches}
\begin{proof}
Use exactly the generated fold of the selected strictly positive form
declaration; its family-field computation passes the function of
recursive values described in the body. This is an explicit premise
about that source declaration, not a theorem granting recursion on
every arbitrary syntax. In the paper's PRED-omega extension, Set
domains can form products into Set, Type and Prop. The target bool
is in Set and the target proposition-code type Prop is in Type, so
both selected folds are licensed, as is elimination of the bool data
discriminator to proposition codes.

For \(\mathcal P\), the five atomic branches are the abstractions
\(\lambda n:\mathsf{nat}.\lambda m:\mathsf{nat}.R(n,m)\),
where R runs through the five Prop relations in the table. Each has
type nat-to-nat-to-Prop. The negation branch is
\(\lambda p:\mathsf{form}.\lambda u:\mathsf{Prop}.\neg u\),
of type form-to-Prop-to-Prop. The three binary branches are
\(\lambda p:\mathsf{form}.\lambda u:\mathsf{Prop}.
\lambda q:\mathsf{form}.\lambda v:\mathsf{Prop}.u\wedge v\),
and the same term with disjunction or implication: each has the
required four-argument branch type. For the two quantifier branches,
in the context \(b:\mathsf{nat},h:\mathsf{nat}\to\mathsf{form},
r:\mathsf{nat}\to\mathsf{Prop}\),
\(\Pi x:\mathsf{nat}.(x<b)\to r x\) and
\(\exists x:\mathsf{nat}.(x<b)\wedge r x\) are propositions.
Three annotated abstractions therefore give precisely the required
quantifier branch type. This lists all eleven branches, including
the unused constructor fields p,q,h that the fold interface supplies.

For comp the atomic branches are the five binary bool comparison
terms already typed. The negation branch is
\(\lambda p:\mathsf{form}.\lambda u:\mathsf{bool}.
\mathsf{b\_not}\ u\). The binary branches abstract p,u,q,v
at form,bool,form,bool and apply the appropriate boolean operation
to u and v. The all branch is
\(\lambda b:\mathsf{nat}.\lambda h:\mathsf{nat}\to\mathsf{form}.
\lambda r:\mathsf{nat}\to\mathsf{bool}.\mathsf{b\_all}\ b\ r\),
and ex uses b\_ex. The bounded program typing just proved shows
these results are bool. This again supplies all eleven typed branches.
Applying the respective folds yields \(\mathcal P:\mathsf{form}\to
\mathsf{Prop}\) and comp:form-to-bool. Their generated iota
computations substitute the corresponding recursive results, giving
exactly every equation in the body table.

Now P(p), the conjunction of the two implications in the adequacy
assertion, is Prop for p:form. At a family constructor its generated
induction premise \(\Pi x:\mathsf{nat}.P(h x)\) is also Prop.
Thus all motives and family premises of form induction are admitted.
This argument does not infer an arbitrary higher-order function's
behavior from parametricity or assume extensional equality of functions.
\end{proof}
\end{annexitem}

\begin{annexitem}{annex:p12.har.prcompiler-ok}{thm:p12-har-prcompiler-ok}{The full eleven-branch universal internal proof}
\begin{proof}
Apply the generated form induction at the well-typed motive
\(P(p):=\mathcal P p\leftrightarrow
\mathsf{istrue}(\mathsf{comp}\ p)\). Denote its two induction
maps at p by \(u_p:\mathcal P p\to\mathsf{istrue}(\mathsf{comp}\ p)\)
and \(v_p\) in the reverse direction.
At each of the five atomic constructors the interpreter computations
identify the goal with the corresponding atomic theorem. Those five
proofs give exactly the five atomic branch functions, generalized
over their n and m fields.

For not p, a proof \(a:\mathcal P p\to\mathsf{FF}\) gives
\(\lambda z. a(v_p z):\neg\mathsf{istrue}(\mathsf{comp}\ p)\).
The reverse boolean-negation equivalence maps this to the required
istrue result. Conversely, its forward boolean equivalence gives
\(b:\neg\mathsf{istrue}(\mathsf{comp}\ p)\); return
\(\lambda z.b(u_p z)\) as a proof of \(\neg\mathcal P p\).
This states explicitly why both directions of the induction assertion
are needed even to construct its forward negation branch.

For and p q, in the forward direction project its two source
components, apply \(u_p,u_q\), pair them and apply the reverse
boolean-conjunction equivalence. In reverse, apply the forward
boolean equivalence, project the resulting pair, apply \(v_p,v_q\)
and pair the source results. For or p q, eliminate the source
disjunction into the boolean istrue proposition, use \(u_p\) with
left injection or \(u_q\) with right injection, and apply the reverse
boolean-disjunction equivalence. In reverse use its forward map,
eliminate the resulting disjunction, and use \(v_p\) or \(v_q\)
with the corresponding source injection. These give both binary
connective branches without extracting any data.

For implication p q, source \(a:\mathcal P p\to\mathcal P q\)
gives the actual map \(\lambda z.u_q(a(v_p z))\) between the two
istrue propositions. The reverse boolean-implication equivalence
returns the required istrue value. In reverse, the forward boolean
equivalence gives \(b:\mathsf{istrue}(\mathsf{comp}\ p)\to
\mathsf{istrue}(\mathsf{comp}\ q)\); return
\(\lambda z.v_q(b(u_p z))\). No classical implication
decision principle is being assumed: its boolean equivalence was
proved by data-constructor cases.

For all b h, the family induction hypothesis supplies \(u_{h x}\)
and \(v_{h x}\) for every x, in the context x:nat, without a
restriction to already enumerated x. Given
\(a:\Pi x.(x<b)\to\mathcal P(h x)\), form
\(\lambda x.\lambda d.u_{h x}(a x d)\).
The reverse b\_all theorem turns this bounded family into the
required istrue result, after the comp constructor computation.
Conversely the forward b\_all theorem gives a bounded family c of
istrue proofs; return \(\lambda x.\lambda d.v_{h x}(c x d)\).
This is the complete universal constructor branch, including the
dependence of the family on its freely chosen x.

For ex b h, eliminate a source witness x and its pair of x<b and
\(\mathcal P(h x)\), into the istrue proposition. Build the same
bounded witness with value proof \(u_{h x}\) applied to the second
component; the reverse b\_ex theorem supplies istrue. Conversely
the forward b\_ex theorem gives a witness x, its bound, and its
istrue value; eliminate that existential only into the source
existential proposition, and build the same witness using \(v_{h x}\).
This completes the eleventh branch.

The generated form induction applied to these eleven branch proof
functions is the actual proof term ok. The preceding descriptions
specify its lambda, conjunction, disjunction, existential and transport
operations, with the only recursions supplied by the already explicit
induction proofs. Generalization over all free declarations, or
substitution by any terms of their declared types, preserves the
typing of every branch. Consequently open parameter contexts and
arbitrary typed higher-order h families are covered; neither closure
nor normalization of their nat fields was used.
\end{proof}
\end{annexitem}

\begin{annexitem}{annex:p12.har.prcompiler-extraction}{thm:p12-har-prcompiler-extraction}{Universal characteristic instances and the exact conversion step}
\begin{proof}
In a context of k nat variables, apply the typed family p to exactly
these arguments to obtain a form value. Apply comp to obtain bool,
and the actual code selector to obtain nat. Abstract over these k
variables to give \(K_p:T_k\). At k=0 the input p is a form value
and the result is the nat datum code(comp p), without an empty-tuple
or unit assumption. The range proof for code, instantiated at
comp(p\(\vec x\)), gives its O/one disjunction for every tuple.
Composition of ok's two maps with the two code-bridge maps gives
\(\mathcal P(p\vec x)\leftrightarrow K_p\vec x=\mathsf{one}\).

For completeness also construct the negative characteristic maps.
Let A be \(\mathcal P(p\vec x)\) and B be
\(\mathsf{istrue}(\mathsf{comp}(p\vec x))\). The adequacy pair
u:A-to-B and v:B-to-A gives \(\neg A\to\neg B\) by
\(a\mapsto\lambda b.a(v b)\), and \(\neg B\to\neg A\) by
\(c\mapsto\lambda a.c(u a)\). Compose these with the negative
code bridge to obtain the asserted iff with zero. This is a
constructive contravariant composition, not an invocation of DN.

For a closed input whose comp result has a specified finite
beta/delta/iota reduction to true, conversion identifies its istrue
proposition with TT. Thus trivial has the required istrue type and
the reverse component of ok p applied to trivial has type
\(\mathcal P p\). This proof construction is conditional on that
actual contraction, not on the bare assertion that a user intended
the formula to be true. Its same universal adequacy remains valid
for an opaque closed or open input which does not reduce.

Finally finite arithmetic expressions built from the compiled
prescriptions have typed nat outputs by the compiler typing theorem.
Put those outputs into the five atomic form constructors, use the
four connective constructors for the matching finite formula
connectives, and for a bounded quantifier abstract its bound variable
as \(\lambda x:\mathsf{nat}.\text{the recursively constructed form}\).
This external construction of a specified finite formula is by
structural recursion on that formula. Constructor equations show
that its \(\mathcal P\) interpretation is the intended formula
on those \emph{same typed arithmetic terms}; the numerical
representation theorem connects numeral valuations to their
independently specified arithmetic values. It supplies instances of
the internal generic program, not an assumed internal inverse on
arbitrary propositions. In particular neither that inverse nor a
classifier of every host nat function follows from ok.
\end{proof}
\end{annexitem}
